\begin{abstract}
Traffic-accident anticipation is increasingly scored by composite metrics that add an earliness term --- time-to-accident at a fixed risk threshold --- to discrimination terms such as average precision and area under the ROC curve. The composition fails structurally rather than incidentally: the discrimination terms are bounded on the unit interval, the earliness terms scale with clip length, and the earliness terms have a trivial global maximiser --- a constant risk score above the threshold attains, on every clip, the largest value that clip admits. We measure the consequence on a live benchmark of 1,417 dashcam clips and thirteen teams, whose labels are withheld and whose metric weights are undisclosed. A constant risk score of 0.51, which reads no pixels, scores 2.35291 on the private leaderboard: it matches the eighth-placed team at the precision the leaderboard displays, exceeds five of thirteen entries including our own, and falls 0.14874 short of the winner. Twenty-five such submissions, designed so that the undisclosed terms cancel when two are differenced, then decompose the metric from outside the competition and without any label: discrimination contributes 0.35105, the two timing terms 2.00186 --- 5.70 times as much --- and the score falls at almost exactly one sixtieth of a point per frame of delay, confirmed by a one-frame experiment. The same probes place the withheld accident-onset distribution, recovering a mean of 120.1 frames in a 150-frame clip, and establish a second result: four curves that cross the threshold at the same frame and never fall below it afterwards score 0.091 apart, a spread the published definition cannot produce. The benchmark's stated formula does not reproduce its own scorer. We close with a protocol whose central requirement is that every anticipation result be reported alongside a video-blind control, which costs one submission. Our own ninth-placed entry, which scores below the constant, is the case study.
\end{abstract}

\noindent\textbf{Keywords:} Accident anticipation, evaluation methodology, benchmark design, video-blind baselines, composite metrics, time-to-accident, dashcam video analysis, intelligent transportation systems.

\begingroup\renewcommand{\thefootnote}{}%
\footnotetext{An AI assistant (Anthropic's Claude) was used in preparing this submission: to draft and edit prose, to generate the video-blind submission files and the scripts that analyse the returned scores, and to convert the manuscript to this template. The research question, the experimental design, the claims and the conclusions are the authors' own, and the authors take full responsibility for all content.}\endgroup

\codeavail
